% ===========================================================================
% QUICK EXAMPLES
% ===========================================================================
%
% The examples below show how to use the specification template.
% They can be copied and adapted directly into the main document.
%
% ---------------------------------------------------------------------------
% Functional Requirement
% ---------------------------------------------------------------------------
%
% \fr{User Authentication}{
%     \item The system shall allow users to authenticate using
%     their username and password.
%
%     \item The system shall reject invalid credentials.
%
%     \item The system shall lock the account after three
%     consecutive failed attempts.
% }
%
% ---------------------------------------------------------------------------
% Non-functional Requirement
% ---------------------------------------------------------------------------
%
% \nfr{Authentication Response Time}{
%     \item The system shall authenticate users within 2 seconds
%     under normal operating conditions.
% }
%
% ---------------------------------------------------------------------------
% Use Case
% ---------------------------------------------------------------------------
%
% \begin{usecase}{Perform Login}
%
%     \ucdescription{
%         Allows a user to authenticate into the system using
%         their credentials.
%     }
%
%     \ucactors{
%         \item User
%     }
%
%     \ucpreconditions{
%         \item The user must be registered in the system.
%         \item The system must be available.
%     }
%
%     \ucmainflow{
%         \item The user enters their username and password.
%         \item The system validates the provided credentials.
%         \item The system creates an authenticated session.
%         \item The system grants access to the application.
%     }
%
%     \ucalternativeflow{
%         \item If the credentials are invalid, the system informs
%         the user that the username or password is incorrect.
%
%         \item If the account is locked, the system denies access.
%     }
%
%     \ucbusinessrules{
%         \item The provided password must match the registered password.
%         \item The account is locked after three consecutive failed attempts.
%     }
%
%     \ucpostconditions{
%         \item The user is authenticated.
%         \item An authenticated session exists for the user.
%     }

%     \ucrelatedrequirements{
%         \item FR01
%     }
%
% \end{usecase}
%
% ---------------------------------------------------------------------------
% Operation Contract
% ---------------------------------------------------------------------------
%
% \begin{operationcontract}{Create Order}
%
%     \ocoperation{
%         createOrder(customer)
%     }
%
%     \ocdescription{
%         Creates a new order for the specified customer.
%     }
%
%     \ocpreconditions{
%         \item The customer must exist.
%         \item The customer must be active.
%     }
%
%     \ocpostconditions{
%         \item A new order is created.
%         \item The order is associated with the customer.
%         \item The order status is set to \texttt{OPEN}.
%     }
%
%     \ocrelatedusecases{
%         \item \ref{uc:place-order}
%     }
%
% \end{operationcontract}
%
% ---------------------------------------------------------------------------
% Referencing
% ---------------------------------------------------------------------------
%
% Requirements:
%
% \ref{fr:user-authentication}
% \ref{nfr:authentication-response-time}
%
% Use cases:
%
% \ref{uc:perform-login}
%
% Operation contracts:
%
% \ref{oc:create-order}
%
% The identifiers (FR01, NFR01, UC01, OC01, ...) are generated
% automatically, and labels are generated from the item names.
%
% ===========================================================================


% ===========================================================================
% PACKAGES
% ===========================================================================
\usepackage[utf8]{inputenc}
\usepackage[a4paper,margin=2cm]{geometry}
\usepackage{enumitem}
\usepackage{titlesec}
\usepackage{hyperref}
\usepackage{titling}
\usepackage{graphicx}
\usepackage[most]{tcolorbox}

% ---------------------------------------------------------------------------
% Customizations
% ---------------------------------------------------------------------------

% Subsection formatting
\titleformat{\subsection}
  {\normalfont\large\bfseries}
  {\thesubsection}
  {1em}
  {\underline}

% Itemize formatting
\setlist[itemize,1]{label=$\bullet$}
\setlist[itemize,2]{label=$\circ$}
\setlist[itemize,3]{label=$\diamond$}

% ---------------------------------------------------------------------------
% Title
% ---------------------------------------------------------------------------

\setlength{\droptitle}{-2cm}

\pretitle{\begin{flushleft}\LARGE\bfseries}
\posttitle{\end{flushleft}}

\preauthor{\begin{flushleft}\large}
\postauthor{\end{flushleft}}

\predate{\begin{flushleft}\large}
\postdate{\end{flushleft} \vspace{-1cm}}

% ---------------------------------------------------------------------------
% Requirements
% ---------------------------------------------------------------------------

\newcounter{functionalrequirement}
\newcounter{nonfunctionalrequirement}

% Functional requirement identifier: FR01, FR02, ...
\renewcommand{\thefunctionalrequirement}{%
    FR\ifnum\value{functionalrequirement}<10 0\fi\arabic{functionalrequirement}%
}

% Non-functional requirement identifier: NFR01, NFR02, ...
\renewcommand{\thenonfunctionalrequirement}{%
    NFR\ifnum\value{nonfunctionalrequirement}<10 0\fi\arabic{nonfunctionalrequirement}%
}

\ExplSyntaxOn

% Convert requirement name into a stable label.
\cs_new_protected:Npn \requirement_slug:n #1
{
    % Convert to lowercase.
    \tl_set:Nx \l_tmpa_tl
    {
        \str_lowercase:n { #1 }
    }

    \str_replace_all:Nnn \l_tmpa_tl { á } { a }
    \str_replace_all:Nnn \l_tmpa_tl { à } { a }
    \str_replace_all:Nnn \l_tmpa_tl { â } { a }
    \str_replace_all:Nnn \l_tmpa_tl { ã } { a }
    \str_replace_all:Nnn \l_tmpa_tl { ä } { a }

    \str_replace_all:Nnn \l_tmpa_tl { é } { e }
    \str_replace_all:Nnn \l_tmpa_tl { è } { e }
    \str_replace_all:Nnn \l_tmpa_tl { ê } { e }
    \str_replace_all:Nnn \l_tmpa_tl { ë } { e }

    \str_replace_all:Nnn \l_tmpa_tl { í } { i }
    \str_replace_all:Nnn \l_tmpa_tl { ì } { i }
    \str_replace_all:Nnn \l_tmpa_tl { î } { i }
    \str_replace_all:Nnn \l_tmpa_tl { ï } { i }

    \str_replace_all:Nnn \l_tmpa_tl { ó } { o }
    \str_replace_all:Nnn \l_tmpa_tl { ò } { o }
    \str_replace_all:Nnn \l_tmpa_tl { ô } { o }
    \str_replace_all:Nnn \l_tmpa_tl { õ } { o }
    \str_replace_all:Nnn \l_tmpa_tl { ö } { o }

    \str_replace_all:Nnn \l_tmpa_tl { ú } { u }
    \str_replace_all:Nnn \l_tmpa_tl { ù } { u }
    \str_replace_all:Nnn \l_tmpa_tl { û } { u }
    \str_replace_all:Nnn \l_tmpa_tl { ü } { u }

    \str_replace_all:Nnn \l_tmpa_tl { ç } { c }

    % Replace every sequence of non-alphanumeric characters with '-'.
    \regex_replace_all:nnN
        { [^a-z0-9]+ }
        { - }
        \l_tmpa_tl

    % Remove '-' from the beginning.
    \regex_replace_once:nnN
        { \A-+ }
        { }
        \l_tmpa_tl

    % Remove '-' from the end.
    \regex_replace_once:nnN
        { -+\Z }
        { }
        \l_tmpa_tl
}

% Functional requirement
\NewDocumentCommand{\fr}{m +m}
{
    \refstepcounter{functionalrequirement}
    \noindent
    \textbf{\thefunctionalrequirement{} - #1}
    \requirement_slug:n { #1 }
    \exp_args:Nx \label { fr:\l_tmpa_tl }

    \begin{itemize}[noitemsep, topsep=2pt, leftmargin=2em]
        #2
    \end{itemize}

    \par\medskip
}

% Non-functional requirement
\NewDocumentCommand{\nfr}{m +m}
{
    \refstepcounter{nonfunctionalrequirement}
    \noindent
    \textbf{\thenonfunctionalrequirement{} - #1}
    \requirement_slug:n { #1 }
    \exp_args:Nx \label { nfr:\l_tmpa_tl }

    \begin{itemize}[noitemsep, topsep=2pt, leftmargin=2em]
        #2
    \end{itemize}

    \par\medskip
}

\ExplSyntaxOff

% ---------------------------------------------------------------------------
% Use Cases
% ---------------------------------------------------------------------------

\newcounter{usecase}

% Use case identifier: UC01, UC02, ...
\renewcommand{\theusecase}{%
    UC\ifnum\value{usecase}<10 0\fi\arabic{usecase}%
}

% Use case box
\newtcolorbox{usecasebox}{
    enhanced,
    breakable,
    lines before break=15,
    colback=white,
    colframe=black!40,
    boxrule=0.5pt,
    arc=2pt,
    left=8pt,
    right=8pt,
    top=6pt,
    bottom=6pt,
    before skip=8pt,
    after skip=12pt,
    segmentation style={
        line width=0.5pt,
        solid
    }
}

\ExplSyntaxOn

% Use case environment
\NewDocumentEnvironment{usecase}{m}
{
    \refstepcounter{usecase}

    % Generate stable label from the use case name.
    \requirement_slug:n { #1 }
    \exp_args:Nx \label { uc:\l_tmpa_tl }

    % Compact title, similar to requirements.
    \begin{usecasebox}
    \noindent
    \textbf{\theusecase{} - #1}
    \par\smallskip
    \tcblower
}
{
    \end{usecasebox}
}

\ExplSyntaxOff

% Use Case Fields ----------------------------------------
\newcommand{\ucdescription}[1]{
    \noindent
    \textbf{Description}
    \par
    \leftskip=1em
    #1
    \par
    \leftskip=0pt
    \medskip
}

\newcommand{\ucactors}[1]{
    \noindent
    \textbf{Actors}
    \begin{itemize}[noitemsep, topsep=2pt, leftmargin=2em]
        #1
    \end{itemize}
}

\newcommand{\ucrelatedrequirements}[1]{
    \noindent
    \textbf{Related requirements}
    \begin{itemize}[noitemsep, topsep=2pt, leftmargin=2em]
        #1
    \end{itemize}
}

\newcommand{\ucpreconditions}[1]{
    \noindent
    \textbf{Preconditions}
    \begin{itemize}[noitemsep, topsep=2pt, leftmargin=2em]
        #1
    \end{itemize}
}

\newcommand{\ucmainflow}[1]{
    \noindent
    \textbf{Main Flow}
    \begin{enumerate}[noitemsep, topsep=2pt, leftmargin=2.4em]
        #1
    \end{enumerate}
}

\newcommand{\ucalternativeflow}[1]{
    \noindent
    \textbf{Alternative Flows}
    \begin{itemize}[noitemsep, topsep=2pt, leftmargin=2em]
        #1
    \end{itemize}
}

\newcommand{\ucbusinessrules}[1]{
    \noindent
    \textbf{Business Rules}
    \begin{itemize}[noitemsep, topsep=2pt, leftmargin=2em]
        #1
    \end{itemize}
}

\newcommand{\ucpostconditions}[1]{
    \noindent
    \textbf{Postconditions}
    \begin{itemize}[noitemsep, topsep=2pt, leftmargin=2em]
        #1
    \end{itemize}
}

% ---------------------------------------------------------------------------
% Operation Contracts
% ---------------------------------------------------------------------------

\newcounter{operationcontract}

% Operation contract identifier: OC01, OC02, ...
\renewcommand{\theoperationcontract}{%
    OC\ifnum\value{operationcontract}<10 0\fi\arabic{operationcontract}%
}

% Operation contract box
\newtcolorbox{operationcontractbox}{
    enhanced,
    breakable,
    lines before break=15,
    colback=white,
    colframe=black!40,
    boxrule=0.5pt,
    arc=2pt,
    left=8pt,
    right=8pt,
    top=6pt,
    bottom=6pt,
    before skip=8pt,
    after skip=12pt,
    segmentation style={
        line width=0.3pt,
        draw=black!30
    }
}

\ExplSyntaxOn

% Operation contract environment
\NewDocumentEnvironment{operationcontract}{m}
{
    \refstepcounter{operationcontract}

    % Generate stable label from the operation contract name.
    \requirement_slug:n { #1 }
    \exp_args:Nx \label { oc:\l_tmpa_tl }

    \begin{operationcontractbox}
        \noindent
        \textbf{\theoperationcontract{} - #1}

        \tcblower
}
{
    \end{operationcontractbox}
}

\ExplSyntaxOff

% Operation Contract Fields ------------------------------------------
\newcommand{\ocoperation}[1]{
    \noindent
    \textbf{Operation}
    \par
    \leftskip=1em
    \texttt{#1}
    \par
    \leftskip=0pt
    \medskip
}

\newcommand{\ocdescription}[1]{
    \noindent
    \textbf{Description}
    \par
    \leftskip=1em
    #1
    \par
    \leftskip=0pt
    \medskip
}

\newcommand{\ocpreconditions}[1]{
    \noindent
    \textbf{Preconditions}
    \begin{itemize}[noitemsep, topsep=2pt, leftmargin=2em]
        #1
    \end{itemize}
}

\newcommand{\ocpostconditions}[1]{
    \noindent
    \textbf{Postconditions}
    \begin{itemize}[noitemsep, topsep=2pt, leftmargin=2em]
        #1
    \end{itemize}
}

\newcommand{\ocrelatedusecases}[1]{
    \noindent
    \textbf{Related Use Cases}
    \begin{itemize}[noitemsep, topsep=2pt, leftmargin=2em]
        #1
    \end{itemize}
}
